\begin{pdSolution}{swk:ex:swk-organs-inmemory}
In memory the substrate still serializes, so resource exclusion, identity and presence, communication, obligation checks, policy checks, and self-attestation keep their transient meaning while the process lives. Everything the subsystems promise across a crash becomes vacuous: durable continuity, crash recovery, historical audit, and outcome evidence; durable identity collapses to an ephemeral handle. The decomposition survives; the cross-crash guarantees do not.
\end{pdSolution}
\begin{pdSolution}{swk:ex:swk-organs-home-table}
The communication subsystem's home table is \texttt{messages}; it owes the layer above an ordered, cursor-addressable, durable envelope with explicit at-least-once semantics. Equally valid: \texttt{claims} for the resource subsystem (one compatible holder per conflict domain), or \texttt{commitments} for the obligation subsystem (no \emph{done} without a current structured oracle receipt).
\end{pdSolution}
\begin{pdSolution}{swk:ex:swk-organs-cut}
A defensible nine-way cut: storage and durability; transaction serialization; identity and presence; resources and leases; message transport; stigmergy and projections; policy and capability; commitments and outcomes; attestation and recovery. It separates proof obligations that leak across the seven, above all transport semantics from message semantics and policy prevention from post-commit detection. It hides the simplifying fact that all local state shares one transaction boundary, so any such cut must retain a cross-cutting invariant: one ledger, one effect mediator.
\end{pdSolution}
\begin{pdSolution}{swk:ex:swk-schema-rename}
Module A renames \texttt{users.name} to \texttt{display\_name}; module B, loaded later under older code, sees no \texttt{name} column and lazily adds it back. Writes now split between two columns according to module version and load order, and neither backfill is authoritative. Every destructive migration has this shape; only additive tables and columns are safe under any-module-any-order.
\end{pdSolution}
\begin{pdSolution}{swk:ex:swk-claim-routing}
Command line, authenticated Unix-socket request, daemon request queue, the sole read/write connection, one transaction, WAL commit. The discipline holds by routing at two points, the client-to-API boundary and the daemon's sole-connection boundary; SQLite's file lock is one backstop at write time. The primary story is routing, not a database rule that forbids other connections.
\end{pdSolution}
\begin{pdSolution}{swk:ex:swk-op7-sequencer}
Expand, backfill, dual-read-or-write, contract keeps mixed versions safe for a while, but the contract step and the dependency order are global facts. Store a migration ledger with a schema epoch, each module's compatibility range, its dependencies, and a backfill watermark. Idempotent self-initialization can remain for additive tables and columns; safe rename, backfill, and down-migration cannot preserve unrestricted any-module-any-order. A \texttt{user\_version} sequencer, or an equivalent ordered capability ledger, is unavoidable.
\end{pdSolution}
\begin{pdSolution}{swk:ex:swk-fault-classes}
(a) I1a-survivable. (b) An orderly shutdown flushes, so a genuinely clean reboot is not the I1b case; an abrupt reset during the reboot is I1b-exposed. (c) I1b-exposed. (d) I1a-survivable after commit. None of the four lets an uncommitted transaction survive.
\end{pdSolution}
\begin{pdSolution}{swk:ex:swk-checkpoint-window}
\par\smallskip
\begin{center}
\begin{tikzpicture}[x=1mm,y=1mm]
  % Timeline axis
  \draw[->,line width=0.8pt,hhink] (0,8) -- (82,8) node[right,font=\footnotesize\selectfont] {$t$};
  % Tick marks and labels
  \draw[line width=0.9pt,hhink] (15,6.5) -- (15,9.5) node[below=11pt,font=\scriptsize\selectfont,align=center] {\textbf{Commit $t_0$}\\(I1a-safe)};
  \draw[line width=0.9pt,hhink] (65,6.5) -- (65,9.5) node[below=11pt,font=\scriptsize\selectfont,align=center] {\textbf{Checkpoint $t_0+\delta$}\\(\texttt{fsync} / I1b-safe)};
  % Vulnerability window shading
  \fill[hhamber!18,draw=hhamber,line width=0.6pt,dashed] (15,10) rectangle (65,18);
  \node[font=\scriptsize\bfseries,text=hhamber!90!hhink] at (40,14) {Vulnerability Window $\delta$ (I1b-exposed)};
  \node[font=\tiny,text=hhgray] at (40,21) {Unsynced WAL frames held in volatile OS page cache};
\end{tikzpicture}
\end{center}
\par\smallskip
The commit is power-loss exposed from $t_0$ until the WAL frames that contain it are synced. A page-count threshold bounds accumulated WAL frames, not elapsed time: at a low or zero write rate it may not fire for arbitrarily long, so it gives no wall-clock bound. Only with an independently guaranteed write-rate floor $w_{\min}$ pages per second and threshold $P$ does $\delta \lesssim P/w_{\min}$ follow, and this chapter assumes no such floor.
\end{pdSolution}
\begin{pdSolution}{swk:ex:swk-op10-interface}
Make durability a typed argument of the transaction, \texttt{PROCESS\_CRASH} or \texttt{POWER\_LOSS}, never an adjective in prose. Route \texttt{POWER\_LOSS} operations through a connection configured \texttt{synchronous=FULL} before the transaction opens and verify the commit before returning; if SQLite settings are connection-scoped, a separate fully-synced financial ledger is the cleanest isolation. Ordinary claims stay on \texttt{NORMAL}. Hard-reset fault tests must exercise the stronger path. The guard against silent defaulting is the type: a write path that names no durability class does not compile, and a forced checkpoint bounds recovery and log size but is not the per-write primitive.
\end{pdSolution}
\begin{pdSolution}{swk:ex:swk-release-idempotent}
A release with no matching (key, holder) row is a successful no-op so that retries are idempotent: this is I3. It must never delete another holder's row. If it raised, a client that timed out on its first release could not distinguish success from failure, and at-least-once handling would break.
\end{pdSolution}
\begin{pdSolution}{swk:ex:swk-acquire-race}
\par\smallskip
\begin{center}
\begin{tikzpicture}[x=1mm,y=1mm]
  % Vertical life-lines
  \draw[line width=0.6pt,pdink!40] (10,32) -- (10,0);
  \draw[line width=0.6pt,pdink!40] (45,32) -- (45,0);
  \draw[line width=0.6pt,pdink!40] (80,32) -- (80,0);

  \node[font=\footnotesize\bfseries,text=pdcobalt] at (10,35) {Agent A};
  \node[font=\footnotesize\bfseries,text=pdink] at (45,35) {Kernel WAL};
  \node[font=\footnotesize\bfseries,text=pdamber] at (80,35) {Agent B};

  % Concurrent acquires
  \draw[->,line width=0.7pt,pdcobalt] (10,29) -- (43,27) node[midway,above,font=\tiny] {\texttt{acquire(k)}};
  \draw[->,line width=0.7pt,pdamber] (80,29) -- (47,27) node[midway,above,font=\tiny] {\texttt{acquire(k)}};

  % First insert succeeds
  \node[draw=pdcobalt,fill=pdcobalt!10,rounded corners=1pt,font=\tiny,inner sep=2pt,align=center] at (45,22)
    {\textbf{INSERT (k, A)}: commits (1st)};
  \draw[->,line width=0.7pt,pdcobalt] (43,18) -- (10,16) node[midway,below,font=\tiny] {\textbf{GRANTED}};

  % Second insert conflicts
  \node[draw=pdamber,fill=pdamber!12,rounded corners=1pt,font=\tiny,inner sep=2pt,align=center] at (45,12)
    {\textbf{INSERT (k, B)}: \textsc{UniqueViolation}};

  % Catch block & select holder
  \node[draw=pdink!70,fill=pdink!5,rounded corners=1pt,font=\tiny,inner sep=2pt,align=center] at (45,5)
    {\texttt{SELECT holder} $\to$ \textbf{A}};
  \draw[->,line width=0.7pt,pdamber] (47,1) -- (80,-1) node[midway,below,font=\tiny] {\textbf{DENIED(A)}};
\end{tikzpicture}
\end{center}
\par\smallskip
The daemon serializes the two requests. The first \texttt{INSERT}\allowbreak(key, holder) commits. The second hits the unique-key violation, selects the extant row inside the catch path, and returns \texttt{DENIED}\allowbreak\texttt{(holder}\allowbreak\texttt{=first)}: that select is where the loser learns the name. For range and hierarchy claims the trace is safe only after both requests map to the same canonical conflict key or the overlap test runs in one immediate transaction.
\end{pdSolution}
\begin{pdSolution}{swk:ex:swk-op1-ticket}
Add a durable FIFO ticket queue per conflict domain, a maximum lease $L$, a monotonically increasing fencing epoch, and grant-to-head on release, on acquire, and on lazy expiry. No background scheduler is required for safety, but liveness requires some live caller or tick to trigger the transitions. Under daemon availability, eventual requests, and non-renewable bounded leases, a requester with $q$ predecessors waits at most about $(q+1)L$ plus processing delay; without those assumptions no bounded wait exists. Every external effect must validate the fencing epoch, or an expired holder can still act.
\end{pdSolution}
\begin{pdSolution}{swk:ex:swk-decay-honest}
If the evaporation tick stalls for an hour, tick-only storage still reports yesterday's marker weight as today's. Read-time evaluation computes $w\,r^{\Delta t}$ from the persisted timestamp and never presents the stale stored number as current.
\end{pdSolution}
\begin{pdSolution}{swk:ex:swk-redelivery}
The transport delivers the request twice. The consumer applies one semantic state transition and recognizes the duplicate as transport nondeterminism; diagnostics may count the redelivery, the operator's view shows one act. An envelope idempotency key plus a unique (consumer, key) receipt lets the consumer store ``effect applied'' atomically with the effect, making retry safety explicit rather than conventional. Exactly-once is still not guaranteed for an unbrokered external effect.
\end{pdSolution}
\begin{pdSolution}{swk:ex:swk-op8-decay}
For each marker kind, log creation, reads, the actions it influenced, its supersession, and the time at which operators label it stale. Fit a survival curve and choose the half-life that minimizes $C_{\text{stale}}\cdot\text{false-live} + C_{\text{lost}}\cdot\text{false-expired}$, validated on held-out projects. Claims and presence should decay far faster than decisions or learned skills. Publish the confidence and re-estimate under drift.
\end{pdSolution}
\begin{pdSolution}{swk:ex:swk-monitor-rules-precommit}
Note-count monotonicity can be regimented by append-only storage or a trigger that rejects deletes and regressive updates. Capability attenuation and lock-owner release can be pre-checked only if authenticated identity, the capability order, and every relevant effect path pass through the mediator. Duplicate keys and boot admission are already pre-effect. Heartbeat freshness cannot be rejected at the heartbeat insert: staleness is an absence observed only after time passes, and under asynchrony a slow actor is indistinguishable from a failed one. Out-of-band filesystem and network effects remain detect-only until mediated.
\end{pdSolution}
\begin{pdSolution}{swk:ex:swk-empty-oracle}
\texttt{close(done, oracle\_ref="")} reaches the oracle-reference check, returns \texttt{REFUSED}, leaves the commitment open, and records the denied authority act. It never reaches classification, current-state verification, or the \texttt{DONE} transition: the empty reference fails the finite-vocabulary test before any oracle is consulted.
\end{pdSolution}
\begin{pdSolution}{swk:ex:swk-classify-rules}
Note monotonicity: the prohibited event is a write the daemon mediates and the trigger (the previous count) is committed state; top-right, regimentable, and append-only storage is the supervisor. Capability escalation: regimentable only if every effect path and the identity behind it pass through the mediator; an escalation exercised through an unmediated channel is an uncontrollable event, which moves the rule to the left column until the channel is owned. Lock-owner validity: the release is mediated and the owner is committed state, top-right, provided the actor identity is authenticated (I12); with a self-asserted identity the trigger is not witnessed and the rule drops to the bottom-right cell.
\end{pdSolution}
\begin{pdSolution}{swk:ex:swk-classify-two}
The first gates a controllable event on a controllable trigger: both \texttt{api\_call} and \texttt{spawn\_child} lie in $\Sigma_c$, the policy never disables an uncontrollable event, and it is regimentable. The second must disable \texttt{internal\_plan}, an uncontrollable event, in its post-trigger state, and the plant enables it there; the history \texttt{fs\_write} followed by \texttt{internal\_plan} witnesses the failure of controllability, so the rule is detect-only.
\end{pdSolution}
\begin{pdSolution}{swk:ex:swk-synthesis}
$P_4$ is a two-state taint automaton whose trigger is \texttt{git\_push}. In the product with $P_2$ (no \texttt{git\_push}, ever) the trigger is disabled in every reachable state, so $P_4$'s tainted state is unreachable in the closed loop and the disablement map is unchanged: \texttt{ok} disables $\{\texttt{net\_egress}, \texttt{git\_push}\}$ and \texttt{tainted} disables $\{\texttt{net\_egress}, \texttt{git\_push}, \texttt{fs\_write}\}$. Controllability still holds, since no reachable state disables an uncontrollable event. The lesson is the one synthesis is for: $P_4$ is vacuous under $P_2$, and a hand-written engine would have carried the dead rule forever.
\end{pdSolution}
\begin{pdSolution}{swk:ex:swk-op5-oracle}
A protected-run receipt: tree hash, a suite hash the principal controls, runner or image hash, command, timestamp and nonce, exit code, and a signed result. Closure verifies every binding and the freshness. This resists free-text and stale-result manipulation; it is not universally Goodhart-proof, because the suite may be incomplete or poisoned. The honest theorem is syntactic provenance and non-replay, not semantic correctness.
\end{pdSolution}
\begin{pdSolution}{swk:ex:swk-notes-vs-work}
Notes preserve explicit goals, rationale, observations, and artifact pointers. They cannot preserve unflushed edits, process memory, open handles, provider caches or hidden state, or a latent next action the model never externalized. A successor can understand the story and still be unable to resume the exact computation.
\end{pdSolution}
\begin{pdSolution}{swk:ex:swk-note-monotonicity}
The append commits. A later delete or regressive count also commits and enters the operation log. Only then does the subscriber observe the count drop and raise or compensate. The violation is durable at the second commit; the monitor proves detection, not prevention.
\end{pdSolution}
\begin{pdSolution}{swk:ex:swk-op4-capsule}
A content-addressed task capsule sealed at a safe tool or message boundary: base tree and worktree blobs, environment and tool manifests, open claims with epochs, commitments and acceptance criteria, exact tool input and output, pending external operations with their idempotency keys, decisions and hypotheses, a provenance-preserving transcript or compaction, and external references. Flush the ledger and quiesce mediated effects before sealing. This restores observable task state and evidence; it cannot recover hidden activations, unexposed reasoning, provider-internal caches, or unrecorded intent. Arbitrary-provider recovery is semantic continuation, not bit-identical resurrection.
\end{pdSolution}
\begin{pdSolution}{swk:ex:swk-crimes}
\texttt{delegate} from the initial state with a child scope not contained in the root grant's scope $\{a,b\}$, for instance $\{a,b,c\}$: with the attenuation guard off the grant is created and I4 fails immediately. No other guard can be broken in one step, because effects need the executing phase, settlement needs the verified phase, and a double settlement needs two operations.
\end{pdSolution}
\begin{pdSolution}{swk:ex:swk-seventh-invariant}
One candidate: a contested work unit never settles (\texttt{settle} is guarded on the verified phase, but the transition from contested back to a settleable phase is unconstrained). Add the guard, rerun the search, and confirm the state count and that the mutation suite finds a trace with the guard off; if the checker finds none, the invariant was already implied and does not earn its row.
\end{pdSolution}
\begin{pdSolution}{swk:ex:swk-linearizability-assumptions}
One daemon connection handles every claim and lock read and write; no \texttt{read\_uncommitted}; checks and mutations share a transaction; respond only after commit; no out-of-band writer; a crash-free observed history. An external writer violates the single-writer and no-bypass conditions. A read-only connection needs its own API and snapshot analysis. Synchronized clocks are unnecessary: non-overlapping invocation/response intervals constrain the order.
\end{pdSolution}
\begin{pdSolution}{swk:ex:swk-linearization}
A invokes \texttt{claim(k)} and commits at $c_1$. B and C overlap; B's denied read commits at $c_2$, A releases at $c_3$, C's acquire commits at $c_4$. One valid order is A-acquire, B-denied, A-release, C-acquire. If B and C overlap around the release, either may linearize first, as long as each result matches the state at its point and non-overlapping operations keep their real-time order.
\end{pdSolution}
\begin{pdSolution}{swk:ex:swk-op3-differential}
Generate identical state-machine traces against both bindings: schema initialization and migrations, acquire, reclaim, release, and expiry, overlapping calls, transactions and rollback, busy and lock behavior, encodings, pragmas, crash and reopen, fault injection. After every step compare returned values, errors, serialized rows, schema, and recovery state, including the deployed binary in CI. A seeded fuzz misses divergences that depend on wall-clock timing, on a pragma the fuzz never sets, or on a crash at a point the fuzz never interrupts; targeted fault injection and pragma sweeps catch those. No finite suite makes I11 a theorem: that needs a formal refinement or equivalence proof, or the same verified binding on both sides. Label the suite conformance evidence.
\end{pdSolution}
\begin{pdSolution}{swk:ex:swk-partial-failure}
The port stays claimed and a live or ghost session row exists, but no commitment says what work owns them. Other agents see a busy port and a presence they cannot reconcile; the failed starter may retry into its own leftovers.
\end{pdSolution}
\begin{pdSolution}{swk:ex:swk-op11-begin-work}
Expose one \texttt{begin\_work} endpoint that opens \texttt{BEGIN IMMEDIATE}, validates every input, inserts the claim, session, commitment, and outbox rows, and commits once; on any error it rolls back all of them. A saga would release the claim and mark or delete the session when the commitment insert failed, and each compensation can itself fail or crash between steps. All three rows live in one SQLite file, so the local transaction is both simpler and stronger.
\end{pdSolution}
\begin{pdSolution}{swk:ex:swk-hash-chain-adversary}
A later same-user tamperer of the database file, or a remote synchronizer that cannot rewrite an externally anchored root. The same-user process is out of scope at the substrate; the adversary becomes real at the cross-machine economy layer. Without an independent anchor or verifier, the writer can rewrite the log and recompute the chain.
\end{pdSolution}
\begin{pdSolution}{swk:ex:swk-forged-release}
The release path compares the supplied actor identity to the lock owner. If a legacy endpoint accepts a self-asserted identity, the attacker submits the victim's string, the equality check passes, and the owner-scoped delete executes. The lock rule is correct relative to its input; I12 failed to authenticate the input. Complete credential and peer binding must precede every security-relevant write.
\end{pdSolution}
\begin{pdSolution}{swk:ex:swk-op9-hardening}
The minimum useful wall is a daemon under a separate UID (or a VM), a daemon-owned database and key, a Unix-socket ACL plus kernel peer credentials, and brokered git, filesystem, and network effects. An OS keychain protects a key at rest but not arbitrary signing requests from an equally authorized same-user process; sealed memory does not stop same-UID tracing or injection without an OS isolation policy; external root anchoring detects a later log rewrite but does not prevent it. No trusted platform module is required for strong local separation: separate privilege and complete mediation are. After the sandbox the same-user adversary moves to the mediated-effects row, where side channels the sandbox does not enumerate remain in scope.
\end{pdSolution}
